\documentclass{situationist-preprint}

\title{Moral Luck and Responsibility}
\msid{ETH-13-a}
\msfamily{ETH-13}
\msdiscipline{Normative and Applied Ethics}
\mstype{Principal manuscript}
\mskeywords{moral luck; control principle; resultant luck; constitutive luck; culpability; reasons-responsiveness; agent-regret}

\begin{document}
\maketitle

\begin{abstract}
Two drivers drive equally carelessly; a child runs into the road in front of one of them and is killed. Ordinary judgement holds the unlucky driver responsible for a death and the lucky one for no such thing, though the difference between them lay beyond their control. Nagel argued that this conflicts with a principle we also accept, that agents are morally assessable only for what is under their control, and that applying the principle consistently would leave nothing for which anyone is responsible. We argue that moral luck is real and that it does not undermine responsibility. We distinguish the degree of an agent's culpability, which depends only on the quality of will she exercised in acts under her control, from the scope of her responsibility, the things she is responsible for, which extends to the outcomes in which her controlled acts issue. The control principle is true of culpability and false of scope. Since responsibility requires control over one's acts, not over their outcomes, their circumstances or one's own constitution, the luck that affects scope leaves the conditions of responsibility in place.
\end{abstract}

\section{The problem}

\citet{Williams1976} and \citet{Nagel1976} gave the problem of moral luck its name and its modern form. Nagel distinguished four ways in which the moral assessment of an agent seems to depend on factors beyond her control: the way her actions turn out (\emph{resultant} luck), the circumstances in which she finds herself (\emph{circumstantial} luck), the kind of person she is (\emph{constitutive} luck), and the way she is determined by antecedent circumstances (\emph{causal} luck). Against these he set what has come to be called the control principle: we are morally assessable only to the extent that what we are assessed for depends on factors under our control. The principle has the authority of Kant, who held that a good will would shine by its own light, like a jewel, even if it accomplished nothing \citep[4:394]{Kant1785}. Nagel argued that if the principle is applied consistently, it erodes the domain of responsibility until nothing is left, since every act depends on factors beyond the agent's control; and that we cannot give it up, since it is part of the concept of responsibility. The result is a paradox.

Three responses have emerged. The first denies that there is moral luck. On this view, differences in the treatment of lucky and unlucky agents reflect differences in our evidence about them, or differences in what is appropriate to do to them, but not in their moral standing \citep{Richards1986}; the lucky driver is exactly as blameworthy as the unlucky one, and is even, as \citet{Zimmerman1987} argued, responsible to the same degree for what he would have done. \citet{EnochMarmor2007} have defended the denial systematically, arguing that each kind of moral luck either disappears on analysis or is not luck in the relevant sense. The second response accepts Nagel's argument and concludes that, because luck is pervasive, no one is ever morally responsible \citep{Levy2011}. The third accepts moral luck and denies that it undermines responsibility. \citet{Moore1997} defended resultant luck in the criminal law; \citet{Hanna2014} and \citet{Hartman2017} have argued that luck often affects how praiseworthy or blameworthy agents are.

We defend a version of the third response. Moral luck is real, and it does not undermine responsibility. Our version differs from the main accepting views in where it locates luck. It does not say that the unlucky driver is more culpable; it says that he is responsible for more.

\section{Framework}

\begin{definition}[Moral luck]\label{d:luck}
An agent is subject to \emph{moral luck} when some correct moral assessment of her depends on a factor that was not under her control.
\end{definition}

\begin{definition}[Control]\label{d:control}
An agent has \emph{control} over an act when the act issues from a mechanism of practical reasoning that is her own and that responds to reasons: had there been sufficient reason to act otherwise, the mechanism would have recognized it and she would have acted otherwise in a suitable range of cases.
\end{definition}

Definition~\ref{d:control} is the account of free will defended in ACT-01, which follows \citet{FischerRavizza1998}. It is control over acts. Nobody has control in this sense over the outcomes of her acts once they leave her hands, over the circumstances that present her with choices, or over the antecedent processes that made her mechanism what it is.

\begin{definition}[Culpability and scope]\label{d:two}
An agent's \emph{culpability} for an act is the degree to which she is blameworthy for it, as fixed by the quality of will the act expresses: the regard or disregard for others it manifests, given what she knew and could reasonably have known. The \emph{scope} of her responsibility is the set of things she is responsible for: her acts and their outcomes, each of which may be attributed to her as her doing.
\end{definition}

The distinction answers two different questions that the reactive attitudes \citep{Strawson1962} and the practices built on them ask about an agent. \emph{How bad was what she did, as an expression of her will?} is a question about culpability, and it governs the intensity of the blame she deserves. \emph{What did she do?} is a question about scope, and it governs what she must answer for: to whom she owes apology, what she owes in reparation, what she can rightly regret as her own doing. With the distinction in hand, the control principle splits in two.

\begin{principle}[Control of culpability]\label{p:culp}
An agent's culpability depends only on factors under her control together with the evidence available to her.
\end{principle}

\begin{principle}[Control of scope]\label{p:scope}
An outcome is within the scope of an agent's responsibility only if it is under her control.
\end{principle}

\begin{principle}[Transmission]\label{p:transmission}
If an agent controls an act, and an outcome results from the act through the realization of a risk that the act culpably created and that she knew or could reasonably have known it created, the outcome is within the scope of her responsibility.
\end{principle}

Principle~\ref{p:culp} is Kant's jewel, restricted to culpability, and we accept it. Principle~\ref{p:scope} is the extension of the control principle to scope, and we reject it. Principle~\ref{p:transmission} is the ground for rejecting it. It is close to the principle by which the law attributes results to those who negligently or recklessly cause them, and its epistemic clause matches the condition defended in ETH-23, that agents are blameworthy only for what they could reasonably have known.

\section{The argument}

\begin{theorem}[Resultant moral luck]\label{t:real}
Moral luck is real.
\end{theorem}

\begin{proof}
Consider the two drivers. Each controls his act of driving carelessly, in the sense of Definition~\ref{d:control}, and each could reasonably have known that careless driving creates a risk of killing a pedestrian. In the first case the risk is realized: a child runs into the road and is killed. By Principle~\ref{p:transmission}, the child's death is within the scope of the first driver's responsibility. In the second case no child runs into the road, and no death is within the scope of the second driver's responsibility. The assessment that the first driver is responsible for a death, which is correct, depends on whether a child ran into the road, which was not under his control. By Definition~\ref{d:luck}, the first driver is subject to moral luck.
\end{proof}

Theorem~\ref{t:real} is established for resultant luck, which is the clearest case; it is enough for the thesis that moral luck is real. The proof shows that Principle~\ref{p:scope} is false, given Principle~\ref{p:transmission}. It does not show that the two drivers differ in culpability, and by Principle~\ref{p:culp} they do not: their wills were the same, and so was their evidence.

\begin{proposition}[Equal culpability, unequal scope]\label{p:equal}
The two drivers are equally culpable, and the unlucky driver is responsible for more.
\end{proposition}

\begin{proof}
Culpability is fixed by quality of will and evidence (Definition~\ref{d:two}), which are the same in both cases, and by Principle~\ref{p:culp} nothing else affects it. The difference in scope is the conclusion of the proof of Theorem~\ref{t:real}.
\end{proof}

Proposition~\ref{p:equal} explains the data that have made the problem hard. It explains why a spectator who knows both cases thinks the drivers equally bad, and why the law nonetheless distinguishes careless driving from causing death by careless driving. It explains why the unlucky driver owes the child's parents an apology and reparation that the lucky driver owes no one: these are owed for what one did, and what he did includes the death. It also explains why the unlucky driver's agent-regret, in Williams's sense, is fitting and the lucky driver's would not be: agent-regret is directed at one's own doing, and the death is part of his. The intensity of the blame the drivers deserve follows culpability and is the same; what each must answer for follows scope and differs.

\begin{theorem}[Responsibility is not undermined]\label{t:notundermined}
Moral luck does not undermine responsibility: each kind of luck leaves in place the conditions under which agents are responsible.
\end{theorem}

\begin{proof}
The condition of responsibility is control over one's acts in the sense of Definition~\ref{d:control}; responsibility for outcomes is transmitted from it by Principle~\ref{p:transmission}. We consider the four kinds of luck in turn.

\emph{Resultant} luck concerns the outcomes of controlled acts. It affects scope, as Theorem~\ref{t:real} shows, but it affects neither the agent's control over her act nor her culpability, by Principle~\ref{p:culp}.

\emph{Circumstantial} luck concerns which choices the agent faces. Nagel's example is the ordinary citizen of Nazi Germany who would have behaved as badly as many did, had he not emigrated for business reasons in 1930. He is not culpable for the acts he did not perform: culpability attaches to the exercise of will in acts under control, and he exercised none. His counterpart who stayed is culpable for acts he controlled. Circumstantial luck determines which acts there are to assess; it does not remove control over the acts there are.

\emph{Constitutive} luck concerns the agent's character and capacities, which she did not choose. Control in the sense of Definition~\ref{d:control} requires that the mechanism from which she acts be her own and responsive to reasons; it does not require that she have created it. ACT-01 argues that a mechanism can be one's own in the required sense although it was produced by processes one did not control, and the argument applies here.

\emph{Causal} luck concerns the antecedent determination of the agent's acts. ACT-01 argues that reasons-responsive control is compatible with determinism and with indeterminism, and ACT-15 that indeterministic luck in the moment of choice is compatible with responsibility when the agent tries to bring about each of the competing outcomes.

In none of the four cases does luck remove the condition of responsibility.
\end{proof}

Nagel's regress argument depends on Principle~\ref{p:scope}, or on a principle that extends control from acts to their antecedents. If responsibility for $X$ required control over everything on which $X$ depends, then responsibility would require control over the outcomes, circumstances, constitution and causes of one's acts, and nobody would be responsible for anything. Theorem~\ref{t:notundermined} rests on the denial of that requirement. Responsibility requires control over one's acts; the rest is the domain of luck, and it is the domain in which moral luck operates.

\paragraph{Retrospect.} The scope of an agent's responsibility is often settled only after the event, since it depends on how the risks she took are realized. This matches a feature of permissibility. ETH-15 argues that every apparent dilemma has a uniquely permissible option, but one that may be identifiable only in retrospect, because the facts that settle it lie partly in the future. An agent who chooses well on her evidence and turns out to have taken the impermissible option is not more culpable for that, by Principle~\ref{p:culp}; but what she did is what she did, and it falls within the scope of her responsibility. Williams's painter, whose decision to abandon his family is vindicated or condemned by whether his work succeeds, is a case of this kind \citep{Williams1976}. The agent may be unable to know at the time of choice what she will turn out to be responsible for, and this is a form of moral luck that leaves her control, at the time of choice, untouched.

\section{Objections}

\begin{objection}[Epistemic reduction]
We treat the unlucky driver differently only because his carelessness has been revealed by its result. Were we to know that the lucky driver was equally careless, we would assess him in the same way \citep{Richards1986}.
\end{objection}

\begin{reply}
Proposition~\ref{p:equal} agrees with this for culpability, and the reduction is right about that. It is wrong about scope. Suppose we know everything about both drivers. We still judge that the unlucky one killed a child and the lucky one did not; that the unlucky one owes the parents an apology and reparation and the lucky one owes them nothing; that the unlucky one can rightly feel agent-regret for a death and the lucky one cannot. These differences survive full information, so they are not differences in evidence.
\end{reply}

\begin{objection}[Counterfactual responsibility]
\citet{Zimmerman1987} argues that the lucky driver is responsible for what he would have done, and so is responsible to the same degree as the unlucky one. Moral luck vanishes when responsibility is extended to counterfactual outcomes.
\end{objection}

\begin{reply}
What Zimmerman calls counterfactual responsibility is culpability under another name: it measures the quality of the will and the risk it accepted, and Principle~\ref{p:culp} already makes it equal. But nobody owes reparation for a death that did not occur, and nobody can apologize to parents whose child is alive. The extension of responsibility to counterfactual outcomes makes the degree of responsibility equal by abandoning its scope, and so leaves unexplained the very practices of answerability that scope governs.
\end{reply}

\begin{objection}[Hard luck]
\citet{Levy2011} argues that present luck, luck in the agent's state of mind at the moment of decision, is pervasive and undermines control itself. If so, luck does not merely affect scope; it removes the condition of responsibility.
\end{objection}

\begin{reply}
Present luck affects which considerations happen to occur to the agent at the moment of choice. Control in the sense of Definition~\ref{d:control} is a property of the mechanism, and it is assessed across a range of cases: a mechanism is reasons-responsive if it would recognize sufficient reasons and act on them in a suitable range of circumstances, including the variations that present luck introduces. An agent whose mechanism is responsive across that range controls her act even if, on a given occasion, which considerations surfaced was not up to her. ACT-01 develops this reply, and ACT-15 its application to indeterministic choice.
\end{reply}

\begin{objection}[Luck in scope is not moral luck]
If luck affects only the set of outcomes for which an agent is responsible, and not how blameworthy she is, it does not affect her moral standing. The real question was whether luck affects how good or bad an agent is, and on the present account it does not \citep{EnochMarmor2007}.
\end{objection}

\begin{reply}
The moral standing of an agent includes what she has done. That a person has killed a child, even through carelessness no worse than that of thousands of others, is a moral fact about her: it shapes her relations to others, grounds what she owes them, and is a fitting object of her own regret. Definition~\ref{d:luck} counts any correct moral assessment, and the assessment that she is responsible for a death is among them. To restrict moral standing to culpability is to adopt Kant's jewel as the whole of morality, which is the conclusion the objection was meant to support.
\end{reply}

\section{What the result does not establish}

The argument does not show that luck affects culpability. Some defenders of moral luck hold that the unlucky driver is more blameworthy \citep{Hartman2017}; Principle~\ref{p:culp} denies this, and our thesis does not need it. Nor does the argument settle how far Principle~\ref{p:transmission} extends: outcomes that result from a controlled act only through the free intervention of another agent, or through risks no one could reasonably have foreseen, raise questions that the principle leaves open. The companion manuscript ETH-13-b works through the two drivers in detail, and ETH-27 applies the distinction between infringing and wrongdoing to the remorse appropriate to agents who act rightly.

\section{Conclusion}

Moral luck is real: what an agent is responsible for depends on how the risks she takes are realized, which is not under her control, and so do the apology, reparation and regret that are owed for it. It does not undermine responsibility, because responsibility requires control over one's acts, and luck in outcomes, circumstances, constitution and causes leaves that control in place. The control principle is true of culpability and false of scope, and Nagel's paradox arises from failing to distinguish them.

\begin{thebibliography}{99}\small
\bibitem[Enoch and Marmor(2007)]{EnochMarmor2007} Enoch, D. and Marmor, A. (2007). The case against moral luck. \emph{Law and Philosophy}, 26(4), 405--436.
\bibitem[Fischer and Ravizza(1998)]{FischerRavizza1998} Fischer, J.~M. and Ravizza, M. (1998). \emph{Responsibility and Control: A Theory of Moral Responsibility}. Cambridge University Press.
\bibitem[Hanna(2014)]{Hanna2014} Hanna, N. (2014). Moral luck defended. \emph{No\^{u}s}, 48, 683--698.
\bibitem[Hartman(2017)]{Hartman2017} Hartman, R.~J. (2017). \emph{In Defense of Moral Luck: Why Luck Often Affects Praiseworthiness and Blameworthiness}. New York: Routledge.
\bibitem[Kant(1785)]{Kant1785} Kant, I. (1785). \emph{Grundlegung zur Metaphysik der Sitten}. Riga: Hartknoch. Cited by Academy edition page.
\bibitem[Levy(2011)]{Levy2011} Levy, N. (2011). \emph{Hard Luck: How Luck Undermines Free Will and Moral Responsibility}. Oxford University Press.
\bibitem[Moore(1997)]{Moore1997} Moore, M.~S. (1997). \emph{Placing Blame: A Theory of Criminal Law}. Oxford: Clarendon Press.
\bibitem[Nagel(1976)]{Nagel1976} Nagel, T. (1976). Moral luck. \emph{Proceedings of the Aristotelian Society, Supplementary Volumes}, 50, 137--151.
\bibitem[Richards(1986)]{Richards1986} Richards, N. (1986). Luck and desert. \emph{Mind}, 95(378), 198--209.
\bibitem[Strawson(1962)]{Strawson1962} Strawson, P.~F. (1962). Freedom and resentment. \emph{Proceedings of the British Academy}, 48, 1--25.
\bibitem[Williams(1976)]{Williams1976} Williams, B. (1976). Moral luck. \emph{Proceedings of the Aristotelian Society, Supplementary Volumes}, 50, 115--135.
\bibitem[Zimmerman(1987)]{Zimmerman1987} Zimmerman, M.~J. (1987). Luck and moral responsibility. \emph{Ethics}, 97(2), 374--386.
\end{thebibliography}

\end{document}
